% !TeX root = ../../MRG41/MRG41_01_poster.tex
% 顶部组织名保留原有的 4 点字距，并随全文字距一起调整。
\pgfmathsetmacro{\OrganizerLetterSpacing}{4+\GlobalLetterSpacing}
\pdfstringdefDisableCommands{\def\\{ }}
\hypersetup{
    hidelinks,
    pdftitle={\EventName\space\SessionNumber\space—\space\PaperTitle},
    pdfauthor={\OrganizerName},
    pdfsubject={\PresenterName, \PresenterAffiliation; \EventDate; \EventTime}}

% 内部派生尺寸和细节间距；通常无需修改。
\newlength{\ContentWidth}
\setlength{\ContentWidth}{\dimexpr\PosterWidth-2\SideMargin\relax}
\newlength{\PresenterAffiliationGap}
\setlength{\PresenterAffiliationGap}{6mm}
\newlength{\DateWeekdayGap}
\setlength{\DateWeekdayGap}{4mm}
\newlength{\WeekdayTimeGap}
\setlength{\WeekdayTimeGap}{5mm}
\newcommand{\eventPrimary}{\normalfont\type{\EventPrimarySize}\medium\color{Ink}}
\newcommand{\eventSecondary}{\normalfont\type{\EventSecondarySize}\color{Muted}}
\newcommand{\eventLabel}{\normalfont\type{\EventLabelSize}\medium\color{Muted}}

% 按输入顺序编号：1–2 位作者纵排；3 位及以上每行两位。
\ExplSyntaxOn
\seq_new:N \g_mrg_authors_seq
\dim_new:N \l_mrg_author_width_dim
\NewDocumentCommand{\AddPaperAuthor}{m m}
    {\seq_gput_right:Nn \g_mrg_authors_seq {{#1}{#2}}}
\cs_new_protected:Npn \mrg_author_card:nnn #1#2#3
    {
        \begin{minipage}[t]{\l_mrg_author_width_dim}
            \raggedright\parindent=0pt
            {
                \fontsize{\AuthorNameFontSize}{\dimexpr\AuthorNameFontSize pt+3pt\relax}
                \selectfont\medium\color{white}
                \int_eval:n {#1}.\hspace{0.5em}#2\par
            }
            \vspace{1mm}
            {
                \fontsize{\AuthorAffiliationFontSize}{\dimexpr\AuthorAffiliationFontSize pt+3pt\relax}
                \selectfont\color{Pale}#3\par
            }
        \end{minipage}
    }
\cs_new_protected:Npn \mrg_author_item:nn #1#2
    {\mrg_author_card:nnn {#1} #2}
\cs_new_protected:Npn \mrg_author_at:n #1
    {\exp_args:Nne \mrg_author_item:nn {#1}{\seq_item:Nn \g_mrg_authors_seq {#1}}}
\NewDocumentCommand{\RenderPaperAuthors}{}
    {
        \int_compare:nNnTF {\seq_count:N \g_mrg_authors_seq} < {3}
            {
                \dim_set:Nn \l_mrg_author_width_dim {\linewidth}
                \seq_map_indexed_inline:Nn \g_mrg_authors_seq
                    {
                        \noindent\mrg_author_card:nnn {##1} ##2\par
                        \int_compare:nNnT {##1} < {\seq_count:N \g_mrg_authors_seq}
                            {\vspace{\AuthorRowGap}}
                    }
            }
            {
                \dim_set:Nn \l_mrg_author_width_dim {(\linewidth-\AuthorColumnGap)/2}
                \int_step_inline:nnnn {1}{2}{\seq_count:N \g_mrg_authors_seq}
                    {
                        \noindent\hbox to\linewidth
                            {
                                \mrg_author_at:n {##1}\hspace{\AuthorColumnGap}
                                \int_compare:nNnT {##1+1} < {\seq_count:N \g_mrg_authors_seq+1}
                                    {\mrg_author_at:n {##1+1}}
                                \hfil
                            }\par
                        \int_compare:nNnT {##1+1} < {\seq_count:N \g_mrg_authors_seq}
                            {\vspace{\AuthorRowGap}}
                    }
            }
    }
\ExplSyntaxOff
\PaperAuthors

% 先测量内容，再计算模块边界和整张海报高度。
\newsavebox{\ModuleTwoBox}
\newsavebox{\ModuleThreeBox}
\newsavebox{\ModuleFourBox}
\newlength{\HeaderHeight}
\setlength{\HeaderHeight}{48mm}
\newlength{\ModuleTwoHeight}
\newlength{\ModuleThreeHeight}
\newlength{\ModuleFourHeight}
\newlength{\ModuleThreeTop}
\newlength{\ModuleFourTop}
\newlength{\PosterHeight}
\setlength{\parindent}{0pt}

% 内容测量
% 保存盒子之间不能产生普通空格，因此这些行尾的 % 必须保留。
\begin{document}%
%
% 模块 2：论文标题、作者与期刊。
\sbox{\ModuleTwoBox}{%
    \begin{minipage}{\ContentWidth}%
        \parindent=0pt\parskip=0pt\raggedright\color{white}%
        {\fontsize{\PaperTitleFontSize}{\dimexpr\PaperTitleFontSize pt+4pt\relax}%
            \selectfont\medium\PaperTitle\par}%
        \vspace{\ContentGap}%
        {\fontsize{\PaperSubtitleFontSize}{\dimexpr\PaperSubtitleFontSize pt+5pt\relax}%
            \selectfont\PaperSubtitle\par}%
        \vspace{\ContentGap}%
        \RenderPaperAuthors%
        \vspace{\ContentGap}%
        {\fontsize{\MetadataFontSize}{\dimexpr\MetadataFontSize pt+4pt\relax}%
            \selectfont\color{Pale}\PaperYear\hspace{1em}\PaperJournal\par}%
    \end{minipage}%
}%
%
% 模块 3：摘要两端对齐，允许自动断词；关键词共用同级分隔符。
\sbox{\ModuleThreeBox}{%
    \begin{minipage}{\ContentWidth}%
        \parindent=0pt\parskip=0pt\color{Ink}%
        {\fontsize{13}{16}\selectfont\medium\color{Blue}\AbstractLabel\par}%
        \vspace{\ContentGap}%
        {\fontsize{\AbstractFontSize}{\dimexpr\AbstractFontSize pt*13/10\relax}%
            \selectfont\hyphenpenalty=50\tolerance=1000\PaperAbstract\par}%
        \vspace{\ContentGap}%
        {\fontsize{\KeywordFontSize}{\dimexpr\KeywordFontSize pt+4pt\relax}%
            \selectfont\color{Muted}\PaperKeywords\par}%
    \end{minipage}%
}%
%
% 模块 4：各行共用字号与基线；房间号、会议号共用左边界。
\sbox{\ModuleFourBox}{%
    \begin{tikzpicture}[
        every node/.style={anchor=base west, inner sep=0pt, outer sep=0pt}
    ]%
        \pgfmathsetlengthmacro{\EventPitch}{\EventPrimarySize pt+\EventRowGap}%
        \pgfmathsetlengthmacro{\EventSecond}{2*\EventPitch}%
        \pgfmathsetlengthmacro{\EventThird}{3*\EventPitch}%
        \pgfmathsetlengthmacro{\EventCodeX}{\VenueValueX mm-\SideMargin}%
        %
        % 报告人：姓名与单位。
        \node[font=\eventLabel] at (0,0) {\PresenterLabel};%
        \node[font=\eventPrimary] (presenter)
            at (\LabelColumnWidth,0) {\PresenterName};%
        \node[font=\eventSecondary, xshift=\PresenterAffiliationGap]
            at (presenter.base east) {\PresenterAffiliation};%
        %
        % 时间：日期与时间同字号，星期视觉弱化。
        \node[font=\eventLabel] at (0,-\EventPitch) {\ScheduleLabel};%
        \node[font=\eventPrimary] (date)
            at (\LabelColumnWidth,-\EventPitch) {\EventDate};%
        \node[font=\eventSecondary, xshift=\DateWeekdayGap] (weekday)
            at (date.base east) {\EventWeekday};%
        \node[font=\eventPrimary, xshift=\WeekdayTimeGap]
            at (weekday.base east) {\EventTime};%
        %
        % 线下地点。
        \node[font=\eventLabel] at (0,-\EventSecond) {\LocationLabel};%
        \node[font=\eventSecondary]
            at (\LabelColumnWidth,-\EventSecond) {\LocationName};%
        \node[font=\eventPrimary] at (\EventCodeX,-\EventSecond) {\LocationRoom};%
        %
        % 线上地点。
        \node[font=\eventLabel] at (0,-\EventThird) {\OnlineLabel};%
        \node[font=\eventSecondary]
            at (\LabelColumnWidth,-\EventThird) {\OnlinePlatform};%
        \node[font=\eventPrimary] at (\EventCodeX,-\EventThird) {\OnlineMeetingID};%
    \end{tikzpicture}%
}%
%
% 模块边界：实际内容高度加统一留白，整张海报只固定宽度。
\setlength{\ModuleTwoHeight}{%
    \dimexpr\ht\ModuleTwoBox+\dp\ModuleTwoBox+2\ModulePadding\relax}%
\setlength{\ModuleThreeHeight}{%
    \dimexpr\ht\ModuleThreeBox+\dp\ModuleThreeBox+2\ModulePadding\relax}%
\setlength{\ModuleFourHeight}{%
    \dimexpr\ht\ModuleFourBox+\dp\ModuleFourBox+2\ModulePadding\relax}%
\setlength{\ModuleThreeTop}{\dimexpr\HeaderHeight+\ModuleTwoHeight\relax}%
\setlength{\ModuleFourTop}{\dimexpr\ModuleThreeTop+\ModuleThreeHeight\relax}%
\setlength{\PosterHeight}{\dimexpr\ModuleFourTop+\ModuleFourHeight\relax}%
\typeout{MRG WIDTH=\the\PosterWidth; HEIGHT=\the\PosterHeight}%
\typeout{MRG MODULE4 HEIGHT=\the\ModuleFourHeight}%
%
% 最终画布：先绘制底色，再按测量结果放置内容。
\begin{tikzpicture}[
    every node/.style={inner sep=0pt, outer sep=0pt, anchor=base west, text=Ink}
]
    % 模块 1：白底活动标识。
    \path[use as bounding box] (0,0) rectangle (\PosterWidth,-\PosterHeight);
    \fill[white] (0,0) rectangle (\PosterWidth,-\HeaderHeight);
    \node[
        font=\type{12}\medium\addfontfeatures{LetterSpace=\OrganizerLetterSpacing},
        text=Muted
    ] at (\SideMargin,-13mm) {\OrganizerName};
    \node[font=\type{\EventNameFontSize}\medium]
        at (\SideMargin,-29mm) {\EventName};
    \node[font=\type{14}\medium, text=Blue]
        at (\SideMargin,-40mm) {\SessionLabel\space\SessionNumber};
    \node[anchor=north east]
        at (\dimexpr\PosterWidth-\SideMargin\relax,-11mm)
        {\includegraphics[width=\LogoWidth]{\LogoPath}};

    % 模块 2：深蓝论文区。
    \fill[Blue] (0,-\HeaderHeight) rectangle (\PosterWidth,-\ModuleThreeTop);
    \node[anchor=north west, yshift=-\ModulePadding]
        at (\SideMargin,-\HeaderHeight) {\usebox{\ModuleTwoBox}};

    % 模块 3：浅蓝摘要区。
    \fill[Mist] (0,-\ModuleThreeTop) rectangle (\PosterWidth,-\ModuleFourTop);
    \node[anchor=north west, yshift=-\ModulePadding]
        at (\SideMargin,-\ModuleThreeTop) {\usebox{\ModuleThreeBox}};

    % 模块 4：白底活动区。
    \fill[white] (0,-\ModuleFourTop) rectangle (\PosterWidth,-\PosterHeight);
    \node[anchor=north west, yshift=-\ModulePadding]
        at (\SideMargin,-\ModuleFourTop) {\usebox{\ModuleFourBox}};
\end{tikzpicture}
\end{document}
